\section{Introduction}

\label{sec:intro}

The global AI compute market is projected to reach \$150 billion by 2028, driven by exponential growth in model training and inference demand~\cite{mckinseyai2024}. Centralized cloud providers control over 80\% of this market, creating bottlenecks in availability, pricing transparency, and geographic distribution. Decentralized compute networks promise to unlock underutilized GPU capacity worldwide, but face a fundamental challenge: \emph{how to verify that a remote, untrusted provider correctly executed an AI workload without re-running the entire computation?}

\subsection{The Verification Problem}

Traditional blockchain verification requires every validator to re-execute every transaction (full replay). For AI workloads, this is infeasible:

\begin{itemize}[leftmargin=1.1em]
    \item A single LLM inference call may consume 10ms--10s of GPU time.
    \item A training job may run for hours or days across multiple GPUs.
    \item Full re-execution would cost as much as the original computation, eliminating the economic benefit of decentralization.
\end{itemize}

Existing approaches---optimistic rollups (re-execute on dispute), zero-knowledge proofs (prove correct execution), TrueBit-style verification games (random spot checks)---each have limitations when applied to AI:

\begin{table}[H]
\centering
\small
\begin{tabular}{lccl}
\toprule
\textbf{Method} & \textbf{Cost} & \textbf{Latency} & \textbf{AI Suitability} \\
\midrule
Full replay & $O(C)$ & High & Infeasible \\
Optimistic rollup & $O(1)$* & 7 days* & Dispute lag \\
ZK proof & $O(C \log C)$ & Minutes & GPU ZK immature \\
TrueBit & $O(C/k)$ & Low & No quality check \\
\textbf{PoAI} & $O(C \cdot \alpha)$ & Low & \textbf{Purpose-built} \\
\bottomrule
\end{tabular}
\caption{Verification approaches. $C$ = computation cost, $\alpha \approx 0.003$ is the verification sampling rate. * = optimistic case.}
\label{tab:verification}
\end{table}

\subsection{Contributions}

\begin{enumerate}[leftmargin=1.4em]
    \item Proof of AI (PoAI): A verification mechanism combining TEE attestation with statistical sampling (\S\ref{sec:poai}).
    \item Multi-consensus architecture separating fast-path and settlement-path (\S\ref{sec:consensus}).
    \item AI-native virtual machine (AIVM) with tensor-aware opcodes (\S\ref{sec:aivm}).
    \item Formal security analysis and economic incentive proofs (\S\ref{sec:security}).
    \item Testnet evaluation with 100 validators and 500 workers (\S\ref{sec:evaluation}).
\end{enumerate}

